%
%============= CHAPTER BREAK =================%
%
\chapter{Julius Pollux (ca 178--192 CE)}%
\justifying%
\vspace*{3mm}%
%
\noindent{}Julius Pollux was a rhetorician trained in Athens %
by Hadrian of Tyre.\endnote{Chronopoulos \& Spöri 2023:30.} %
Early in the singular reign of Commodus, %
Pollux was “appointed to the Chair of Rhetoric in Athens”.\endnote{\textit{Ibid}.} %
Caught up in the Attic renaissance of his time, %
he spent several years compiling a massive onomastic dictionary, the \textit{Onomasticon}. %
It follows a notional rather than alphabetical structure, %
and at the micro-level “often adopts a format based on parts of speech and %
the semantic relations of synonymy and antonymy”.\endnote{\textit{Ibid}.} \\[1mm] %

\noindent{}\textbf{σκηνοποιός} / \textit{skēnopoios} / “tent-maker” %
occurs only 121 times in the TLG. Its two earliest hits are the canonical Acts of the Apostles %
and the Athenian rhetorician, Julius Pollux. %
The usage in Acts is a memorable and well-known characterization of Paul, %
his hosts Priscilla and Aquila, and their shared vocation, %
set during his stay in Corinth. %

\vspace*{2mm} 
 \begin{quote}καὶ διὰ τὸ ὁμότεχνον εἶναι ἔμενεν παρ’ αὐτοῖς καὶ ἠργάζετο· %
 ἦσαν γὰρ \textbf{σκηνοποιοὶ} τῇ τέχνῃ. \\[1mm] %
 And because of him being the same-craft, he stayed with them and was working. %
 For they were \textbf{tent-makers} by craft. \end{quote} % 
 \vspace*{2mm}

\noindent{}The cognate lemma σκηνοποιία / \textit{skēnopoiia} / “tent-making”, “scene-making”, “theater-building”, %
or even “nest-building” (LSJ) has a much older history, %
first attested in the TLG by the 4th century BCE Aeneas (\textit{Poliorcetica} 8.3), %
and corroborated by roughly contemporaneous Greek inscriptions.\endnote{%
E.g., τὴν σκηνοπο[ιία]ν / \textit{tēn skēnopo[iia]n} in TAM I 65 Lycia, Ctr. — Isinda (Belenli) — late 5th--mid 4th c. BC?. %
τὴν ἄλλην σκηνοποιίαν / \textit{tēn allēn skēnopoiian} in Hatzopoulos, Mac. Inst. II 12 σκηνοποιίαν Mak. (Edonis) — Amphipolis — ca. 200 BC — ISE II 114 — SEG 40.524. %
LSJ, s.v., σκηνοποι-ία, ἡ, provides additional references, noting its usage in regard to the building of %
theaters and nests, both interesting to consider in relation to the passage in Acts 18:1--3.}

The passage in Pollux runs as follows, taken from \textit{Onomasticon} book 7, %
which was composed between 177 and the early 180s CE.\endnote{%
Chronopoulos \& Spöri 2023: 30--31.} %
We omit the embedded scholarly annotations found in the TLG digital text, %
which provides cross-references to Athenian exemplars, %
even though below we discuss the most relevant cross-reference. %
The specific excerpt comes from \textit{Onomasticon} 7.189.\endnote{Bethe 1900.2:104.} %

\vspace*{2mm} 
 \begin{quote}θαυματοποιός θαυματοποιία, θαυμάτων δημιουργοί. \newline
 μηχανοποιός μηχανοποιία μηχανήματα μηχαναί, \newline
 εὐμήχανοι εὐμηχανία, πολυμήχανοι πολυμηχανία, \newline
 ἀμήχανοι ἀμηχανία, βιομήχανοι ὡς Ἀντιφῶν. \newline
 \textls[-15]{τοὺς δὲ μηχανοποιοὺς καὶ \textbf{σκηνοποιοὺς} ἡ παλαιὰ κωμῳδία ὠνόμαζεν.}\\[1mm] %
 Wonder-worker, wonder-working, creators of wonders. \newline
 Machine-maker, machine-making, machines, mechanisms. \newline
 Inventives, inventiveness, poly-inventives, poly-inventiveness. \newline
 Unskilled, unskillfulness, life-skilled as Antiphon. \newline
 The Old Comedy even called \textbf{scene-makers} machine-makers.\end{quote} % 
 \vspace*{2mm}

As Betheʼs shorthand annotation clarifies, %
the last statement makes a meaningful diachronic linguistic distinction. %
This is confirmed for fragment 98 of the “Anonymous” / \textit{Adespota} section and “Ancient” / \textit{Archaias} subsection %
of the \textit{Comicorum atticorum fragmenta} (Kock 1888:417). %
Kock cites back to this same passage in Pollux: %

\vspace*{2mm} 
 \begin{quote}\textbf{σκηνοποιούς} — τοὺς μηχανοποιοὺς ἡ παλαιὰ κωμῳδία ὠνόμαξεν. \\[1mm] %
 \textbf{Scene-makers} — the Old Comedy called them machine-makers. \end{quote} % 
\vspace*{2mm} 

\noindent{}Another ancient, cognate word also played this role, %
a lemma attested by Plutarch (\textit{De musica} 1142B) for Aristoxenus %
(4th c. BCE): %
σκηνικος / \textit{skēnikos} / “theatrical” or “actor” (LSJ), perhaps even “stager”.\endnote{%
Aristoxenus fr. 76; Wehrli 1967:30.}

Various LSJ entries provide helpful supplemental context about the semantic association %
between machine-craft and stage-craft. %
For example, for the term I translated above as “mechanisms” / μηχαναί (s.v., μηχανή, ‑ῆς, ἡ), %
LSJ notes “theatrical machine by which gods, etc., were made to appear in the air”. %

I did not find the main lemma, σκηνοποιός / \textit{skēnopoios} --- the one in Acts and Julius Pollux ---, %
in searches of major databases of Greek papyri and inscriptions. %
Even when Plutarch speaks of the theater and musical production, %
he does not use this specific lemma. %
The cumulative evidence thus points to its emergence around the time of Pollux, %
perhaps first invented, standardized, and/or popularized by him. %

This all lends support to the cumulative argument of the book. %
This portion of Acts likely reflects an historical and literary setting after Pollux, %
i.e., in the 180s CE or afterwards. %
It also shows by contrast that NT lexica and scholarly works that default to “tent-maker” %
as the termʼs meaning have largely missed the lemmaʼs clear theatrical domain and significance. %
Paul, Priscilla, and Aquila in Acts are likely depicted not as mere “tent-makers”, %
but rather as “scene-makers” or “theater-builders” in Corinth. %
With a nod to Shakespeare, we might say that Acts and its characters are together engaged in stagecraft.\endnote{%
On the literary and mythological theatricality of Acts 19 %
(which has the only occurrence of the word “theater” in Acts), %
and for further literature on the theatricality of canonical Acts, see esp. Bilby \& Lefteratou 2022.} %